\chapter{Information, Entropy and Error-Correcting Codes}\label{ch:b3:information-and-codes}

How many bits does a message really need, and how many more must be added so that it survives a noisy line? Both questions have exact answers, given by Claude Shannon in 1948. % ledger: b3:shannon-1948
The first is the entropy of the source: no code can do better on average, and Huffman's code comes within one bit of it. The second is the capacity of the channel: below it, codes can make errors as rare as wanted. This chapter measures both on Brenmoor's library catalogue, then builds the codes engineers use: the \omterm{def:b3:information-and-codes:hamming}{Hamming code}, which corrects any single flipped bit with three check bits for four of data, and the \omterm{def:b3:information-and-codes:cyclic}{cyclic codes} behind every CRC, read again as division of polynomials. Every claim about error rates is checked twice, by computing it and by simulating a seeded channel.

\begin{recall}
The entropy of a finite distribution $p_1, \ldots, p_n$ is $\sum_i p_i \log_2 (1/p_i)$ (Year~2 volume). A prefix code gives each symbol a binary codeword, no codeword a prefix of another, so a stream decodes without separators; Huffman's algorithm builds one of least average length (Year~1 volume). A cyclic redundancy check appends to a frame the remainder of a division, so that the receiver detects most corruptions (Year~2 volume). Matrices, their products and the field with two elements, where $1 + 1 = 0$, are those of the mathematics book.
\end{recall}

\section{Entropy}\label{sec:b3:information-and-codes:entropy}

\begin{definition}[Self-information, entropy of a source]\label{def:b3:information-and-codes:entropy}
A \emph{source} emits symbols of an alphabet independently, each $x$ with probability $p(x)$. The \emph{self-information}\index{self-information} of an outcome of probability $p$ is $\log_2 (1/p)$ bits. The \emph{entropy of a source}\index{entropy of a source} is the expected self-information of a symbol, $H = \sum_x p(x) \log_2 (1/p(x))$ bits per symbol. The \emph{block entropy} $H_k$ is the entropy of blocks of $k$ symbols taken as one symbol; for a source of independent symbols, $H_k = kH$.
\end{definition}

A rare outcome carries much information, a certain one none, and the information of two independent outcomes adds, since $\log_2 (1/(pq)) = \log_2 (1/p) + \log_2 (1/q)$. The source of this chapter is the titles of the 80 books in Brenmoor's library catalogue, in lower case with letters and spaces only.

\omcode{code/bachelor-3/19-information-and-codes/listings/entropy.py}{9}{17}{Entropy, and the frequencies of blocks of symbols.}\label{lst:b3:information-and-codes:entropy}

\omoutput{code/bachelor-3/19-information-and-codes/outputs/entropy.txt}{The catalogue's letters: the space is the commonest symbol, \texttt{w} the rarest; 3.99 bits per symbol, against 4.52 for a uniform source on 23 symbols.}\label{out:b3:information-and-codes:entropy}

\begin{theorem}[Gibbs' inequality; entropy is at most $\log_2 n$]\label{thm:b3:information-and-codes:gibbs}
For distributions $p$ and $q$ on the same $n$ symbols, $\sum_x p(x) \log_2 (1/q(x)) \ge \sum_x p(x) \log_2 (1/p(x))$, with equality only when $p = q$. In particular $H \le \log_2 n$, with equality only for the uniform source.
\end{theorem}

\begin{proof}
Using $\ln t \le t - 1$, with equality only at $t = 1$, the difference is $\sum_x p(x) \log_2 (q(x)/p(x)) \le \frac{1}{\ln 2} \sum_x p(x) \bigl(q(x)/p(x) - 1\bigr) = \frac{1}{\ln 2} \bigl(\sum_x q(x) - 1\bigr) \le 0$, sums over the $x$ with $p(x) > 0$. Take $q$ uniform: $\sum_x p(x) \log_2 n = \log_2 n \ge H$.
\end{proof}

For two symbols with probabilities $p$ and $1 - p$ the entropy is the binary entropy function $h(p) = p \log_2 (1/p) + (1 - p) \log_2 (1/(1 - p))$: zero at $p = 0$ and $p = 1$, and one bit at $p = 1/2$.

\begin{omfigure}
\begin{tikzpicture}
  \begin{axis}[omrecord, width=0.62\linewidth, height=5cm, xmin=0, xmax=1, ymin=0, ymax=1.05, xlabel={$p$}, ylabel={bits}, legend pos=outer north east, legend cell align=left, legend style={font=\small}]
    \addplot[omDef, thick, domain=0.0005:0.9995, samples=200] {-x*log2(x) - (1-x)*log2(1-x)};
    \addlegendentry{$h(p)$}
    \addplot[omThm, thick, dashed, domain=0.0005:0.9995, samples=200] {1 + x*log2(x) + (1-x)*log2(1-x)};
    \addlegendentry{$1 - h(p)$}
  \end{axis}
\end{tikzpicture}
\par\smallskip{\small The binary entropy function and the capacity $1 - h(p)$ of the \omterm{def:b3:information-and-codes:channel}{binary symmetric channel} of \cref{sec:b3:information-and-codes:channels}.}
\end{omfigure}

Joint and conditional entropy, the uncertainty of a pair of variables and of one when the other is known, extend these definitions; they are not needed here.

\section{The source-coding theorem}\label{sec:b3:information-and-codes:source-coding}

\begin{definition}[Kraft inequality, uniquely decodable code]\label{def:b3:information-and-codes:prefix}
A code is \emph{uniquely decodable}\index{uniquely decodable code} when no two different strings of symbols encode to the same bits; prefix codes are. Codeword lengths $\ell_1, \ldots, \ell_n$ satisfy the \emph{Kraft inequality}\index{Kraft inequality} when $\sum_i 2^{-\ell_i} \le 1$.
\end{definition}

\begin{theorem}[Kraft]\label{thm:b3:information-and-codes:kraft}
The lengths of a prefix code satisfy the \omterm{def:b3:information-and-codes:prefix}{Kraft inequality}, and for any lengths that satisfy it there is a prefix code with those lengths.
\end{theorem}

A codeword of length $\ell$ claims the fraction $2^{-\ell}$ of all infinite bit strings, those that begin with it, and in a prefix code the claims are disjoint (\cref{exo:b3:information-and-codes:7} is the full proof). Kraft published the inequality for prefix codes in 1949; Brockway McMillan proved in 1956 that it holds for every \omterm{def:b3:information-and-codes:prefix}{uniquely decodable code}. % ledger: b3:kraft-mcmillan
% ledger: b3:mcmillan-1956
The construction gives the shortest lengths the smallest binary numbers, a \emph{canonical} code.

\omcode{code/bachelor-3/19-information-and-codes/listings/kraft.py}{6}{23}{The Kraft sum, a prefix test, and the canonical code for given lengths.}\label{lst:b3:information-and-codes:kraft}

\omoutput{code/bachelor-3/19-information-and-codes/outputs/kraft.txt}{A Kraft sum of 1 is not enough to be a prefix code; it is enough for some prefix code with those lengths to exist.}\label{out:b3:information-and-codes:kraft}

\begin{theorem}[Source-coding theorem]\label{thm:b3:information-and-codes:source-coding}\index{source-coding theorem}
For a source of entropy $H$, the least expected length $L^*$ of a prefix code for its symbols satisfies $H \le L^* < H + 1$. Coding blocks of $k$ symbols at once gives an expected length per symbol below $H + 1/k$.
\end{theorem}

The lower bound is Gibbs' inequality with $q(x) = 2^{-\ell(x)}$; the upper bound takes $\ell(x) = \lceil \log_2 (1/p(x)) \rceil$, which satisfies Kraft (\cref{exo:b3:information-and-codes:10}). Huffman's code is optimal, so its length is $L^*$. David Huffman published his construction in 1952. % ledger: b3:huffman-1952
Below, the catalogue's letters are treated as a source of independent letters with the catalogue's frequencies, and Huffman codes are built for single letters, pairs and triples.

\omcode{code/bachelor-3/19-information-and-codes/listings/huffman_vs_h.py}{12}{25}{The lengths of a Huffman code: each merge adds one bit to every symbol beneath it.}\label{lst:b3:information-and-codes:huffman}

\omoutput{code/bachelor-3/19-information-and-codes/outputs/huffman_vs_h.txt}{Huffman against the entropy: within 0.04 bits per letter for single letters, and closer with blocks, as the theorem says; the bound $1/k$ is far from reached.}\label{out:b3:information-and-codes:huffman}

So no lossless code for this source averages fewer than 3.99 bits per letter, and a simple one averages 4.03. The 23-symbol uniform code would need 5 bits, or $\log_2 23 \approx 4.52$ with blocks.

\section{Channels}\label{sec:b3:information-and-codes:channels}

\begin{definition}[Binary symmetric channel, channel capacity, repetition code]\label{def:b3:information-and-codes:channel}
The \emph{binary symmetric channel}\index{binary symmetric channel} with crossover probability $p$ transmits bits, flipping each independently with probability $p$. Its \emph{channel capacity}\index{channel capacity} is $C = 1 - h(p)$ bits per use. A code of \emph{rate} $k/n$ sends $k$ bits of message in $n$ uses of the channel. The \emph{repetition code}\index{repetition code} of length $n$ sends each bit $n$ times, rate $1/n$, and decodes by majority.
\end{definition}

\begin{omfigure}
\begin{tikzpicture}[font=\small, >=Stealth]
  \node (a0) at (0, 2.0) {0}; \node (a1) at (0, 0) {1};
  \node (b0) at (4.4, 2.0) {0}; \node (b1) at (4.4, 0) {1};
  \draw[->, thick] (a0) -- node[above] {$1 - p$} (b0);
  \draw[->, thick] (a1) -- node[below] {$1 - p$} (b1);
  \draw[->, omThm] (a0) -- node[pos=0.22, right=2pt] {$p$} (b1);
  \draw[->, omThm] (a1) -- node[pos=0.22, right=2pt] {$p$} (b0);
  \node[anchor=east] at (-0.3, 1.0) {sent};
  \node[anchor=west] at (4.7, 1.0) {received};
\end{tikzpicture}
\par\smallskip{\small The \omterm{def:b3:information-and-codes:channel}{binary symmetric channel}: each bit arrives intact with probability $1 - p$ and flipped with probability $p$.}
\end{omfigure}

\begin{theorem}[Noisy-channel coding theorem]\label{thm:b3:information-and-codes:noisy-channel}\index{noisy-channel coding theorem}
For every rate $R < C$ and every $\varepsilon > 0$ there are codes of rate at least $R$ whose probability of decoding a block wrongly is below $\varepsilon$; above capacity, no such codes exist.
\end{theorem}

Shannon stated it in 1948; the proof, by choosing a code at random, is not given here. % ledger: b3:noisy-channel
It says that noise costs rate, not reliability. The \omterm{def:b3:information-and-codes:channel}{repetition code} shows what a naive code pays.

\omcode{code/bachelor-3/19-information-and-codes/listings/bsc.py}{8}{15}{The channel, and the exact error of majority decoding.}\label{lst:b3:information-and-codes:bsc}

\omoutput{code/bachelor-3/19-information-and-codes/outputs/bsc.txt}{\omterm{def:b3:information-and-codes:channel}{Repetition codes} over a seeded channel, 100,000 bits each: the measured rates agree with the exact ones.}\label{out:b3:information-and-codes:bsc}

\omoutput{code/bachelor-3/19-information-and-codes/outputs/capacity.txt}{Capacity of the \omterm{def:b3:information-and-codes:channel}{binary symmetric channel}.}\label{out:b3:information-and-codes:capacity}

At $p = 0.1$ the capacity is 0.53 bits per use, so codes of rate $1/2$ with errors as rare as wanted exist; the \omterm{def:b3:information-and-codes:channel}{repetition code} of rate $1/5$ still errs on almost one bit in a hundred. Making a \omterm{def:b3:information-and-codes:channel}{repetition code} reliable sends its rate to zero. Better codes spread each message bit over several check bits.

\section{Hamming distance and linear codes}\label{sec:b3:information-and-codes:linear}

\begin{definition}[Hamming distance, minimum distance, error-correcting code]\label{def:b3:information-and-codes:distance}
The \emph{Hamming distance}\index{Hamming distance} between two words of the same length is the number of positions where they differ. The \emph{minimum distance}\index{minimum distance} of a code is the least distance between two of its codewords. A code is an \emph{error-correcting code}\index{error-correcting code} for $t$ errors when nearest-codeword decoding recovers the sent word whenever at most $t$ bits are flipped.
\end{definition}

\begin{theorem}[Detection and correction]\label{thm:b3:information-and-codes:distance}
A code of \omterm{def:b3:information-and-codes:distance}{minimum distance} $d$ detects every pattern of at most $d - 1$ errors and corrects every pattern of at most $\lfloor (d - 1)/2 \rfloor$.
\end{theorem}

The distance is a metric, so the balls of radius $\lfloor (d - 1)/2 \rfloor$ around the codewords are disjoint, and fewer than $d$ flips cannot turn a codeword into another (\cref{exo:b3:information-and-codes:8}).

\begin{definition}[Linear code, generator matrix, parity-check matrix, syndrome]\label{def:b3:information-and-codes:linear}
A \emph{linear code}\index{linear code} of length $n$ and dimension $k$ is a subspace of $\mathbb{F}_2^n$ of dimension $k$, $\mathbb{F}_2$ being the field with two elements: the sum of two codewords, bit by bit modulo 2, is a codeword. A \emph{generator matrix}\index{generator matrix} is a $k \times n$ matrix $G$ whose rows are a basis: the message $m$ encodes to $mG$. A \emph{parity-check matrix}\index{parity-check matrix} is an $(n - k) \times n$ matrix $H$ with $Hc = 0$ exactly for the codewords $c$. The \emph{syndrome}\index{syndrome} of a received word $r$ is $Hr$.
\end{definition}

If $r = c + e$, the \omterm{def:b3:information-and-codes:linear}{syndrome} is $H(c + e) = He$: it depends on the error only, not on the message. And the \omterm{def:b3:information-and-codes:distance}{minimum distance} of a \omterm{def:b3:information-and-codes:linear}{linear code} is the least weight of a nonzero codeword, since $c - c'$ is a codeword of weight equal to their distance.

\begin{definition}[Hamming code, perfect code, sphere-packing bound]\label{def:b3:information-and-codes:hamming}
The \emph{Hamming code}\index{Hamming code} with $r$ check bits has length $n = 2^r - 1$ and a \omterm{def:b3:information-and-codes:linear}{parity-check matrix} whose columns are the $n$ nonzero vectors of $\mathbb{F}_2^r$, column $j$ being $j$ written in binary. The \emph{sphere-packing bound}\index{sphere-packing bound}: a code of $M$ words of length $n$ correcting $t$ errors has $M \sum_{i=0}^{t} \binom{n}{i} \le 2^n$. A code meeting it with equality is a \emph{perfect code}\index{perfect code}.
\end{definition}

The bound holds because the balls of radius $t$ around the codewords are disjoint and each holds $\sum_{i \le t} \binom{n}{i}$ words of $\mathbb{F}_2^n$. For $r = 3$ the \omterm{def:b3:information-and-codes:hamming}{Hamming code} is the $(7, 4)$ code: a single error at position $j$ has \omterm{def:b3:information-and-codes:linear}{syndrome} column $j$, which is $j$ in binary, so the receiver flips bit $j$ back. No two columns are equal and none is zero, so no codeword has weight 1 or 2, and the \omterm{def:b3:information-and-codes:distance}{minimum distance} is 3. A \omterm{def:b3:information-and-codes:linear}{generator matrix}, its rows being the codewords of the four unit messages, is
\[
G = \begin{pmatrix} 1&1&1&0&0&0&0\\ 1&0&0&1&1&0&0\\ 0&1&0&1&0&1&0\\ 1&1&0&1&0&0&1 \end{pmatrix}.
\]
With 16 codewords and balls of $1 + 7 = 8$ words, $16 \times 8 = 128 = 2^7$: the $(7, 4)$ code is perfect, and every word of 7 bits is within distance 1 of exactly one codeword (\cref{exo:b3:information-and-codes:11}). Richard Hamming published these codes in 1950, after relay computers at Bell Labs kept stopping on errors their punched input could detect but not correct. % ledger: b3:hamming-1950
% ledger: b3:hamming-history

\begin{omfigure}
\begin{tikzpicture}[font=\small]
  \draw[omDef, thick, fill=omDef!8, fill opacity=0.6] (0, 0.9) circle (1.5);
  \draw[omThm, thick, fill=omThm!8, fill opacity=0.6] (1.8, 0.9) circle (1.5);
  \draw[black!60, thick, fill=black!5, fill opacity=0.6] (0.9, -0.65) circle (1.5);
  \node at (0.9, 0.4) {$d_7$};
  \node at (0.9, 1.5) {$d_3$};
  \node at (0.05, -0.2) {$d_5$};
  \node at (1.75, -0.2) {$d_6$};
  \node at (-0.7, 1.4) {$p_1$};
  \node at (2.5, 1.4) {$p_2$};
  \node at (0.9, -1.55) {$p_4$};
  \node[anchor=west, align=left] at (4.0, 0.4) {circle $p_1$: positions 1, 3, 5, 7\\circle $p_2$: positions 2, 3, 6, 7\\circle $p_4$: positions 4, 5, 6, 7\\each circle holds an even number of 1s};
\end{tikzpicture}
\par\smallskip{\small The $(7, 4)$ \omterm{def:b3:information-and-codes:hamming}{Hamming code} as three circles. A single flipped bit makes exactly the circles that contain it odd, and the set of odd circles names its position: that set is the \omterm{def:b3:information-and-codes:linear}{syndrome}.}
\end{omfigure}

\omcode{code/bachelor-3/19-information-and-codes/listings/hamming74.py}{7}{30}{Hamming $(7, 4)$: encode, compute the \omterm{def:b3:information-and-codes:linear}{syndrome}, correct.}\label{lst:b3:information-and-codes:hamming}

\omoutput{code/bachelor-3/19-information-and-codes/outputs/hamming74.txt}{One flipped bit at position 5 is found and corrected; the \omterm{def:b3:information-and-codes:distance}{minimum distance} is 3; and a million 4-bit blocks over the channel with $p = 0.01$.}\label{out:b3:information-and-codes:hamming}

At $p = 0.01$ an uncoded block of four bits is wrong with probability $1 - 0.99^4 \approx 3.94\%$. The \omterm{def:b3:information-and-codes:hamming}{Hamming code} fails only when two or more of its seven bits flip, with probability $1 - 0.99^7 - 7 \times 0.01 \times 0.99^6 \approx 0.203\%$, at rate $4/7$; the simulation agrees.

\section{Cyclic codes}\label{sec:b3:information-and-codes:cyclic}

\begin{definition}[Cyclic code, Reed--Solomon code]\label{def:b3:information-and-codes:cyclic}
Read a word $c_0 c_1 \cdots c_{n-1}$ as the polynomial $c_0 + c_1 x + \cdots + c_{n-1} x^{n-1}$ with coefficients in $\mathbb{F}_2$. A \emph{cyclic code}\index{cyclic code} is a \omterm{def:b3:information-and-codes:linear}{linear code} closed under rotating its words by one place; its codewords are the multiples of a \emph{generator polynomial} $g$ that divides $x^n - 1$. A \emph{Reed--Solomon code}\index{Reed--Solomon code} is a cyclic code over a larger finite field, whose symbols are bytes rather than bits.
\end{definition}

The CRC of the Year~2 volume is division by a generator: a message $m$ of degree below $k$ is sent as $m(x) x^r + (m(x) x^r \bmod g(x))$, a multiple of $g$, and the receiver checks that the remainder of what it receives is zero. Over $\mathbb{F}_2$ subtraction is exclusive or, so long division is a loop of shifts and XORs. A \emph{burst} of length $b$ is an error pattern whose first and last flipped bits are $b - 1$ positions apart.

\omcode{code/bachelor-3/19-information-and-codes/listings/crc_poly.py}{7}{18}{Remainder and encoding over $\mathbb{F}_2$, with the generator $x^8 + x^2 + x + 1$.}\label{lst:b3:information-and-codes:crc}

\begin{theorem}[Burst detection]\label{thm:b3:information-and-codes:burst}
If the generator $g$ has degree $r$ and a nonzero constant term, every burst of length at most $r$ is detected.
\end{theorem}

\begin{proof}
A burst of length $b \le r$ starting at position $i$ is $e(x) = x^i b(x)$ with $b(0) = 1$ and $\deg b = b - 1 < r$. It goes undetected only if $g$ divides $x^i b(x)$. Since $g(0) = 1$, $g$ has no factor in common with $x^i$, so $g$ would divide $b$, a nonzero polynomial of smaller degree: impossible.
\end{proof}

The program checks the theorem by brute force on a 32-bit codeword, then counts what longer bursts escape: a burst of length $r + 1 = 9$ is missed only when it equals $g$ shifted, one pattern in $2^7$, and longer bursts one in $2^8$.

\omoutput{code/bachelor-3/19-information-and-codes/outputs/crc_poly.txt}{Every burst of up to 8 bits is caught; longer bursts escape at the rates $2^{-7}$ and $2^{-8}$.}\label{out:b3:information-and-codes:crc}

\omterm{def:b3:information-and-codes:cyclic}{Cyclic codes} are well suited to bursts, the runs of corrupted bits that scratches and electrical noise produce. % ledger: b3:cyclic-burst
\omterm{def:b3:information-and-codes:cyclic}{Reed--Solomon codes}, by Irving Reed and Gustave Solomon in 1960, use the finite fields of the mathematics book's Year~3 to correct whole corrupted bytes. % ledger: b3:reed-solomon-1960
They are in compact discs, the first mass-produced consumer use in 1982, in DVDs and Blu-ray discs, and in QR codes and other two-dimensional bar codes, which they let a scanner read even with part of the code damaged. % ledger: b3:rs-uses
They are told here, not built.

\section{Exercises}\label{sec:b3:information-and-codes:exercises}

\begin{exercise}[$\star$]\label{exo:b3:information-and-codes:1}
Compute the entropy of a fair six-sided die, and of a die that shows 6 with probability $1/2$ and each other face with probability $1/10$. Which is larger, and why must it be?
\end{exercise}

\begin{exercise}[$\star$]\label{exo:b3:information-and-codes:2}
For each code, give its Kraft sum and say whether it is prefix-free: $\{0, 10, 11\}$; $\{1, 01, 001, 000\}$; $\{0, 10, 01\}$; $\{00, 01, 100, 101, 11\}$.
\end{exercise}

\begin{exercise}[$\star$]\label{exo:b3:information-and-codes:3}
Give the \omterm{def:b3:information-and-codes:distance}{Hamming distances} between \texttt{1011010}, \texttt{1001110} and \texttt{0011011}. How many errors could a code containing these three words correct?
\end{exercise}

\begin{exercise}[$\star$]\label{exo:b3:information-and-codes:4}
\code{hamming74.py} receives \texttt{1110100}. Predict the \omterm{def:b3:information-and-codes:linear}{syndrome}, the corrected codeword and the decoded data.
\end{exercise}

\begin{exercise}[$\star\star$]\label{exo:b3:information-and-codes:5}
\code{exercises/ex5_syndrome.py} decodes Hamming $(7, 4)$ with the rows of the \omterm{def:b3:information-and-codes:linear}{parity-check matrix} taken in the opposite order. Which single errors does it still correct, which does it not, and why? Fix it.
\end{exercise}

\begin{exercise}[$\star\star$]\label{exo:b3:information-and-codes:6}
Write \code{encode15} and \code{decode15} for the Hamming $(15, 11)$ code, parity bits at positions 1, 2, 4 and 8. They are tested on every single error of random words.
\end{exercise}

\begin{exercise}[$\star\star$]\label{exo:b3:information-and-codes:7}
Prove \cref{thm:b3:information-and-codes:kraft}.
\end{exercise}

\begin{exercise}[$\star\star$]\label{exo:b3:information-and-codes:8}
Prove \cref{thm:b3:information-and-codes:distance}.
\end{exercise}

\begin{exercise}[$\star\star$]\label{exo:b3:information-and-codes:9}
Compute the probability that the \omterm{def:b3:information-and-codes:channel}{repetition code} of length 3 decodes a bit wrongly at $p = 0.05$ and $p = 0.2$, and check against \cref{out:b3:information-and-codes:bsc}. What is the probability that a 4-bit block sent with it arrives wrong at $p = 0.01$, and at what rate, compared with Hamming $(7, 4)$?
\end{exercise}

\begin{exercise}[$\star\star\star$]\label{exo:b3:information-and-codes:10}
Prove the source-coding theorem: $H \le L^* < H + 1$, and the bound $H + 1/k$ for blocks of $k$ symbols.
\end{exercise}

\begin{exercise}[$\star\star\star$]\label{exo:b3:information-and-codes:11}
Prove that every \omterm{def:b3:information-and-codes:hamming}{Hamming code}, with $r \ge 2$ check bits, has \omterm{def:b3:information-and-codes:distance}{minimum distance} 3 and is perfect.
\end{exercise}

\begin{exercise}[$\star\star\star$]\label{exo:b3:information-and-codes:12}
Prove that a burst of length $r + 1$ escapes the generator $g$ of degree $r$, with $g(0) = 1$, only when it is a shift of $g$, and deduce the fraction $2^{-(r-1)}$ of \cref{out:b3:information-and-codes:crc}.
\end{exercise}

\section{Weekend problem}\label{sec:b3:information-and-codes:problem}

\begin{problem}[{Weekend problem --- Sending the library catalogue over a noisy line}]\label{pb:b3:information-and-codes:1}
The library sends its catalogue of titles to the town hall over an old line that flips each bit with probability $p = 0.01$. This problem compresses the catalogue, adds a \omterm{def:b3:information-and-codes:hamming}{Hamming code}, and computes and simulates what arrives.

\medskip\noindent\textbf{Part I --- The source.}
\begin{enumerate}
  \item The catalogue has 1,437 symbols over 23 letters. How many bits does a fixed-length code need?
  \item What does an entropy of 3.99 bits per symbol mean for every lossless code?
  \item Which symbol carries the most information, and how much?
  \item How many bits does the Huffman code of single letters take for the whole catalogue, on average per letter as in \cref{out:b3:information-and-codes:huffman}?
  \item Why can coding triples of letters not beat 3.99 bits per letter on this source?
  \item The real titles repeat whole words. Would a model of pairs of letters have lower or higher entropy per letter than the independent model? Why?
\end{enumerate}

\medskip\noindent\textbf{Part II --- The line.}
\begin{enumerate}[resume]
  \item What is the capacity of the line?
  \item The compressed catalogue is about 5,800 bits. How many bit errors are expected with no code?
  \item What is the probability that it arrives with no error at all?
  \item Why is a single error in a Huffman-coded stream worse than in a fixed-length one?
  \item How many channel uses at least does the noisy-channel theorem allow for 5,800 bits sent reliably?
  \item What would a \omterm{def:b3:information-and-codes:channel}{repetition code} of length 3 cost, and what bit error would it leave?
\end{enumerate}

\medskip\noindent\textbf{Part III --- Hamming (7,4).}
\begin{enumerate}[resume]
  \item How many 4-bit blocks and channel bits does the coded catalogue use?
  \item Compute the probability that a block decodes wrongly.
  \item How many wrong blocks are expected in the catalogue?
  \item Compare with the uncoded block error of $1 - 0.99^4$.
  \item What happens when two bits of one block flip?
  \item Why can the code not even tell that two bits flipped?
  \item What do \cref{out:b3:information-and-codes:hamming}'s measured rates add?
\end{enumerate}

\medskip\noindent\textbf{Part IV --- The bill.}
\begin{enumerate}[resume]
  \item Give the rate of the coded transmission and compare it with the capacity.
  \item Would the $(15, 11)$ \omterm{def:b3:information-and-codes:hamming}{Hamming code} do better or worse on this line? Compute its block error.
  \item How would a CRC added to each title help, and what can it not do?
  \item A burst of 5 bits hits the line. What does Hamming $(7, 4)$ do, and what does a CRC of degree 8 do?
  \item How could interleaving, sending the first bits of many blocks, then the second bits, turn a burst into single errors?
  \item Sum up: the bits sent, the expected errors, and what the town hall receives.
\end{enumerate}
\end{problem}
